\section*{Chapitre \ref{ch:b2:chemical-potential} --- Potentiel chimique et mélanges}

\begin{solution}{exo:b2:chemical-potential:1}
La tangente en $x_2 = 0.4$ coupe les bords en $\bar V_1 \approx 14$ et $\bar
V_2 \approx 49$ (unités du modèle). Alors $0.6 \times 14 + 0.4 \times 49 = 28.0$,
le volume molaire lu sur la courbe en $x_2 = 0.4$, inférieur aux
$0.6 \times 18 + 0.4 \times 58 = 34.0$ des volumes purs.
\end{solution}

\begin{solution}{exo:b2:chemical-potential:2}
$n(\text{eau}) = 1000/18.015 = \qty{55.5}{mol}$, $x_1 = 55.5/56.5 = 0.982$~;
$p = 0.982 \times 3.17 = \qty{3.11}{kPa}$~: un abaissement de 1.8\,\%.
\end{solution}

\begin{solution}{exo:b2:chemical-potential:3}
$p(\ce{O2}) = 0.2095 \times \qty{1.013e5}{Pa} = \qty{2.12e4}{Pa}$~;
$c = 1.3 \times 10^{-5} \times 2.12 \times 10^4 = \qty{0.28}{mol/m^3} =
\qty{2.8e-4}{mol/L}$, soit $2.8 \times 10^{-4} \times 32.0 \times 10^3 =
\qty{8.8}{mg/L}$.
\end{solution}

\begin{solution}{exo:b2:chemical-potential:4}
$Q = (1.0)^2/[(1.0)(0.010)^3] = 1.0 \times 10^6$~; $RT\ln Q = 2.479 \times 13.82 =
\qty{34.2}{kJ/mol}$~; $\Delta_r G = -32.8 + 34.2 = \qty{+1.4}{kJ/mol} > 0$~: dans ce
mélange si pauvre en dihydrogène, l'ammoniac se décompose.
\end{solution}

\begin{solution}{exo:b2:chemical-potential:5}
Gibbs--Duhem~: $x_1\dd\bar V_1 + x_2\dd\bar V_2 = 0$, donc $\dd\bar V_2 =
-(0.70/0.30) \times (-0.30) = \qty{+0.70}{mL/mol}$.
\end{solution}

\begin{solution}{exo:b2:chemical-potential:6}
$c(\text{particules}) = 2 \times 9.0/58.44 = \qty{0.308}{mol/L} =
\qty{308}{mol/m^3}$~; $\Pi = 308 \times 8.314 \times 310 = \qty{7.9e5}{Pa}$,
environ \qty{7.9}{bar}. L'intérieur d'un globule rouge a la même \omterm{def:b2:chemical-potential:osmotic-pressure}{pression osmotique}~:
aucun flux net d'eau. Dans l'eau pure, l'eau entre dans la cellule, qui gonfle et
éclate.
\end{solution}

\begin{solution}{exo:b2:chemical-potential:7}
$c = \Pi/RT = 825/(8.314 \times 298.15) = \qty{0.333}{mol/m^3}$~; dans
\qty{1.000e-4}{m^3}, $n = \qty{3.33e-5}{mol}$, $M = 2.00/3.33 \times 10^{-5}
= \qty{6.0e4}{g/mol}$ (\qty{60}{kg/mol}). $\Delta T_f = 1.86 \times 3.33
\times 10^{-5}/0.100 = \qty{6.2e-4}{K}$~: non mesurable, alors que la \omterm{def:b2:chemical-potential:osmotic-pressure}{pression
osmotique} correspond à une colonne de \qty{8}{cm} d'eau.
\end{solution}

\begin{solution}{exo:b2:chemical-potential:8}
(a) $I = \qty{1.0e-3}{mol/kg}$, $\log\gamma_\pm = -0.51 \times 0.0316$,
$\gamma_\pm = 0.964$. (b) $I = \frac12(4 \times 1.0 + 1 \times 2.0) \times
10^{-3} = \qty{3.0e-3}{mol/kg}$, $\log\gamma_\pm = -0.51 \times 2 \times 0.0548$,
$\gamma_\pm = 0.879$. (c) $I = \frac12(4 + 4) \times 10^{-3} = \qty{4.0e-3}{mol/kg}$,
$\log\gamma_\pm = -0.51 \times 4 \times 0.0632$, $\gamma_\pm = 0.74$.
\end{solution}

\begin{solution}{exo:b2:chemical-potential:9}
$b = 0.310/1.86 = \qty{0.167}{mol/kg}$~; $n = 0.167 \times 0.0500 =
\qty{8.33e-3}{mol}$~; $M = 1.50/8.33 \times 10^{-3} = \qty{180}{g/mol}$ (un
sucre en C6, par exemple).
\end{solution}

\begin{solution}{exo:b2:chemical-potential:10}
$x_1\,\dd(Ax_2^2) + x_2\,\dd(Ax_1^2) = 2Ax_1x_2(\dd x_2 + \dd x_1) = 0$ puisque
$x_1 + x_2 = 1$. Quand $x_1 \to 0$, $\gamma_1 \to \mathrm e^{A}$ et $p_1 =
\gamma_1x_1p_1^* \approx \mathrm e^{A}p_1^*\,x_1$~: $k_{H,1} = p_1^*\mathrm e^A$.
Pour $A = 1.1$, $\mathrm e^{1.1} = 3.0$~: la pente de Henry vaut trois fois celle de Raoult.
\end{solution}

\begin{solution}{exo:b2:chemical-potential:11}
(a) Un soluté volatil apporte sa propre pression de vapeur~; la vapeur n'est plus
du solvant pur et la démonstration (vapeur pure du solvant en
équilibre avec la solution) tombe en défaut. (b) Si le soluté entre dans le solide, le
\omterm{def:b2:chemical-potential:chemical-potential}{potentiel chimique} du solide est lui aussi abaissé, de $RT\ln x_1^{\text{s}}$. Si la
solution solide a la même composition que le liquide, les deux droites descendent
de la même quantité et la température de congélation ne change pas~: l'abaissement
exige un soluté rejeté par le cristal.
\end{solution}

\begin{solution}{exo:b2:chemical-potential:12}
Cryométrie~: $b_{\min} = 0.01/1.86 = \qty{5.4e-3}{mol/kg}$~; ébulliométrie~:
$0.01/0.513 = \qty{1.9e-2}{mol/kg}$. À $\qty{5.4e-3}{mol/L} =
\qty{5.4}{mol/m^3}$, $\Pi = 5.4 \times 8.314 \times 298 = \qty{1.3e4}{Pa}$,
plus d'un mètre d'eau~; un osmomètre lit quelques pascals. L'osmométrie est de
loin la plus sensible~: c'est pourquoi elle mesure les grandes masses molaires des
polymères, dont les molalités sont minuscules.
\end{solution}

\begin{solution}{pb:b2:chemical-potential:1}
\textbf{1.} $\mu_1 = \mu_1^*(T, p) + RT\ln x_1$.
\textbf{2.} Diol $30/62.07 = \qty{0.483}{mol}$, eau $70/18.015 =
\qty{3.886}{mol}$~; $x_1 = 3.886/4.369 = 0.889$.
\textbf{3.} $p = 0.889 \times 3.17 = \qty{2.82}{kPa}$.
\textbf{4.} $b = 0.483/0.070 = \qty{6.90}{mol/kg}$.
\textbf{5.} Le diol bout à \qty{197}{\celsius}~: sa pression de vapeur à température
ambiante est bien inférieure à celle de l'eau.
\textbf{6.} $\mu_1^{\text{s}}(T_f) = \mu_1^{*\text{l}}(T_f) + RT_f\ln x_1$.
\textbf{7.} $\ln x_1 = -\Delta_{\text{fus}}G(T)/(RT)$ avec
$\Delta_{\text{fus}}G = \mu_1^{*\text{l}} - \mu_1^{\text{s}}$, nul en $T_f^*$~;
la relation de Gibbs--Helmholtz, $\dd(\Delta_{\text{fus}}G/T)/\dd T = -\Delta_{\text{fus}}H/T^2$,
intégrée de $T_f^*$ à $T_f$, donne $\Delta_{\text{fus}}G(T_f)/T_f =
\Delta_{\text{fus}}H(1/T_f - 1/T_f^*)$, d'où la relation.
\textbf{8.} $\ln x_1 \approx -x_2 \approx -bM_1$ et $1/T_f - 1/T_f^* \approx
-\Delta T_f/T_f^{*2}$~: $\Delta T_f = (RT_f^{*2}M_1/\Delta_{\text{fus}}H)\,b$.
\textbf{9.} $333.4 \times 18.015 = \qty{6007}{J/mol}$~; $K_f = 8.314 \times
273.15^2 \times 0.018015/6007 = \qty{1.86}{K.kg/mol}$.
\textbf{10.} Que $\Delta_{\text{fus}}H$ ne varie pas entre $T_f$ et
$T_f^*$ (elle varie, d'après la \omterm{thm:b2:reaction-enthalpy:kirchhoff}{loi de Kirchhoff}, d'environ 2\,\% pour \qty{3}{K} ici).
\textbf{11.} $\Delta T_f = 1.86 \times 6.90 = \qty{12.8}{K}$~: \qty{-12.8}{\celsius}.
\textbf{12.} $1/T_f = 1/273.15 - (8.314/6007)\ln 0.889$~: $T_f = \qty{261.6}{K}$,
\qty{-11.6}{\celsius}.
\textbf{13.} La relation idéale exacte, qui ne suppose pas la solution
diluée (ici 11\,\% des molécules sont du diol). Toutes deux supposent un \omterm{def:b2:chemical-potential:ideal-mixture}{mélange idéal}
et une enthalpie de fusion constante~; l'eau et le diol, liés par liaisons hydrogène,
ne forment pas un \omterm{def:b2:chemical-potential:ideal-mixture}{mélange idéal}~: la température de congélation réelle diffère.
\textbf{14.} De la glace pure~: le diol reste dans le liquide.
\textbf{15.} Retirer de l'eau sous forme de glace abaisse $x_1$ dans le liquide et, d'après la
relation, la température d'équilibre diminue~: le liquide gèle sur un intervalle
de températures, et non en un point.
\textbf{16.} $K_b = 8.314 \times 373.12^2 \times 0.018015/40\,650 =
\qty{0.513}{K.kg/mol}$.
\textbf{17.} $\Delta T_b = 0.513 \times 6.90 = \qty{3.5}{K}$.
\textbf{18.} $1/T = 1/373.12 - (8.314/40\,650)\ln(2.0/1.013)$~: $T = \qty{393.5}{K}$,
environ \qty{120}{\celsius}.
\textbf{19.} Le bouchon~: \qty{20}{K} contre \qty{3.5}{K} pour le diol.
\textbf{20.} $x_1 = \exp[-(6007/8.314)(1/253.15 - 1/273.15)] = 0.811$.
\textbf{21.} $x_2 = 0.189$~: $w = 0.189 \times 62.07/(0.189 \times 62.07 + 0.811
\times 18.015) = 0.44$.
\textbf{22.} Loi des solutions diluées~: $b = 20/1.86 = \qty{10.75}{mol/kg}$, $w = 10.75 \times
62.07/(1000 + 10.75 \times 62.07) = 0.40$. Les linéarisations ($\ln x_1 \approx
-x_2$, $x_2 \approx bM_1$) surestiment l'abaissement par mole de diol à une concentration
aussi élevée~: la loi des solutions diluées demande donc moins de diol.
\textbf{23.} Dans le modèle du \omterm{def:b2:chemical-potential:ideal-mixture}{mélange idéal}, un liquide de refroidissement de fraction massique en
diol $\boldsymbol{w \approx 0.44}$ reste liquide jusqu'à \qty{-20}{\celsius}.
\end{solution}
